% A544-P0511-HEADING
\clearpage
\SourceLeaf{511}%
\renewcommand{\BookRunningTitle}{Zur Berechnung der gemeinen Logarithmen}
\pdfbookmark[0]{Zur Berechnung der gemeinen Logarithmen}{gemeine-logarithmen}
\begin{center}\Large ZUR BERECHNUNG DER GEMEINEN LOGARITHMEN.\end{center}
% A544-P0511-OPENING
Man suche die Logarithmen von
% A544-P0511-DEFINITIONS
\[
\begin{aligned}
\log\frac{1025}{*}&=a\\
\log\frac{1024^2}{*}&=b\\
\log\frac{81^2}{*}&=c\\
\log\frac{125^2}{*}&=d\\
\log\frac{99^2}{*}&=e
\end{aligned}
\]
% A544-P0511-STAR-EXPLANATION
($*$ zeigt einen um $1$ kleinern Nenner als Zähler an) so ist, wovon man sich leicht überzeugen kann:
% A544-P0511-LOG-TWO
\[
\log 2=\frac{14\frac34+\frac54a+2b-\frac34c-2d+e}{49}=f
\]
% A544-P0511-NEXT-PROSE
Noch kann man leicht herleiten
% A544-P0511-OTHER-IDENTITIES
\[
\begin{aligned}
\log 41&=a+12f-2=g\\
\log 3&=\frac{1+c+4f+g}{8}
\end{aligned}
\]
% A544-P0511-TERMINAL
$\log 11\pdot31$ und $\log\frac{11}{7}$ und also auch $\log 7\pdot31$\par\medskip\noindent\rule{.32\textwidth}{.3pt}
